\documentclass[11pt,a4paper]{article}

\usepackage{lineno}
\linenumbers

\usepackage[utf8]{inputenc}
\usepackage[T1]{fontenc}
\usepackage[british]{babel}
\usepackage[margin=2.5cm]{geometry}
\usepackage{parskip}
\usepackage{enumitem}
\usepackage{titlesec}
\usepackage[hidelinks]{hyperref}

\usepackage{draftwatermark}
\SetWatermarkText{DRAFT}
\SetWatermarkScale{1}

% Section spacing tuned for a short terms-of-reference style document
\titlespacing*{\section}{0pt}{1.4em}{0.6em}
\titlespacing*{\subsection}{0pt}{1.0em}{0.4em}

\setlist[itemize]{itemsep=0.2em, topsep=0.3em}

\title{ACTS Project Contact \\ \large Terms of Reference}
\author{}
\date{\today\\v0.1}

\begin{document}

\maketitle
\thispagestyle{empty}

\section*{Purpose}

The Project Contact (PC) is the named entry point and coordinator for the ACTS
project. The role exists to give the project a recognisable face to the outside
world, to convene discussion when one is needed, and to carry administrative threads
that benefit from a single owner --- without concentrating decision-making power.

The PC does not unilaterally direct the technical evolution of the project.
Authority over the codebase remains distributed among maintainers and contributors,
and the project continues to operate by consensus in its developer meetings and on
the project's code review platform. The PC's role is to keep that process running,
not to override it.

A Deputy Project Contact is appointed alongside the PC to share the load, ensure
continuity, and stand in when the PC is unavailable.

\section*{Responsibilities}

The PC is responsible for:
\begin{itemize}
    \item Acting as the external entry point for the project. This includes
          representing ACTS at funding and oversight bodies, and being the default
          contact for parties outside the immediate developer community.
    \item Chairing the Experiment Liaison Group, or delegating the chair to the
          Deputy or another contributor.
    \item Convening discussion when disputes arise that the normal review process
          has not resolved, ensuring that all affected parties are heard, and
          facilitating a return to consensus.
    \item Setting and communicating the broad direction of the project, with
          community input. ``Direction'' here means roadmap framing and
          prioritisation across competing demands, not unilateral technical
          mandates.
    \item Maintaining the contributor list and coordinating citation and authorship
          questions, in line with the project's separately documented policies on
          those topics.
\end{itemize}

The Deputy shares these responsibilities as needed and steps in for the PC when
required. The split of work between PC and Deputy is left to them to organise.

\section*{Out of scope}

\begin{itemize}
    \item Code review and merge decisions on individual pull requests, which
          continue through the regular review workflow and the existing reviewer
          pool.
    \item Chairing of the regular developer meeting, which continues under its
          existing arrangement.
    \item Decisions internal to any experiment using ACTS.
    \item Anything that would require the PC to make commitments on behalf of
          contributors or their host institutions.
\end{itemize}

\section*{Selection and term}

The PC and Deputy positions are self-perpetuating: the incumbent PC identifies a
successor (and likewise for the Deputy) when stepping down, in consultation with
the active developer community. There is no formal confirmation procedure, and no
fixed term --- the expectation is that succession is discussed openly in the
developer meeting and that the choice reflects the sense of the community.

The initial PC is Andreas Salzburger; the initial Deputy is to be named.

\section*{Review}

These terms of reference are intentionally light, in keeping with the project's
preference for small, reversible steps over heavy formalisation. They will be
reviewed at least annually, alongside the ELG terms of reference, and adjusted as
practice settles.

\end{document}